\documentclass[11pt,a4paper]{article}
\usepackage[utf8]{inputenc}
\usepackage[T1]{fontenc}
\usepackage[spanish,es-noshorthands]{babel}
\usepackage[margin=2.5cm]{geometry}
\usepackage{amsmath,amssymb,mathtools}
\usepackage{enumitem}
\usepackage{booktabs}
\usepackage{tikz}
\usepackage{pgfplots}
\pgfplotsset{compat=1.18}

\newcounter{ejnum}
\newenvironment{ejercicio}{\refstepcounter{ejnum}\par\medskip\noindent\textbf{Ejercicio \theejnum.}~}{\par}
\newenvironment{solucion}{\par\noindent\textit{Solución.}~}{\par\vspace{1em}}

\title{\vspace{-1.5cm}Práctica 2: Análisis combinatorio, probabilidad condicional y teorema de Bayes \\[0.3em] \large Unidad 1: Espacios de probabilidad}
\author{Probabilidad y Estadística (5765) \\ Universidad Nacional del Sur}
\date{}

\begin{document}
\maketitle

\begin{ejercicio}
Un local de comidas rápidas ofrece 5 tipos de hamburguesa, 4 tipos de guarnición y 3 bebidas. Un cliente arma un menú eligiendo una opción de cada categoría.
\begin{enumerate}[label=(\alph*)]
\item ¿Cuántos menús distintos puede armar?
\item Si además puede elegir agregar o no un postre (de entre 6 postres disponibles, o ninguno), ¿cuántos menús distintos hay ahora?
\end{enumerate}
\end{ejercicio}
\begin{solucion}
\begin{enumerate}[label=(\alph*)]
\item Por el principio multiplicativo: $5\times 4\times 3 = 60$ menús distintos.
\item Las opciones de postre son 6 postres más la opción ``ningún postre'', es decir 7 alternativas. Entonces:
\[
5\times 4\times 3\times 7 = 420 \text{ menús distintos.}
\]
\end{enumerate}
\end{solucion}

\begin{ejercicio}
En una línea de producción se deben ensamblar 6 componentes electrónicos distintos en una placa, uno a continuación de otro, y el orden de ensamblado afecta el resultado final (tiempos de secado, accesibilidad, etc.).
\begin{enumerate}[label=(\alph*)]
\item ¿De cuántas formas distintas se puede ordenar el ensamblado de los 6 componentes?
\item Si 2 de los 6 componentes deben ir necesariamente uno junto al otro (en cualquier orden entre ellos, pero adyacentes), ¿de cuántas formas se puede realizar el ensamblado?
\end{enumerate}
\end{ejercicio}
\begin{solucion}
\begin{enumerate}[label=(\alph*)]
\item Es una permutación simple de 6 objetos distintos: $6! = 720$.
\item ``Pegamos'' los dos componentes que deben ir juntos, tratándolos como un único bloque. Ahora hay 5 unidades a ordenar (el bloque + los otros 4 componentes): $5!=120$ formas. Dentro del bloque, los 2 componentes se pueden ordenar de $2!=2$ formas. Por el principio multiplicativo:
\[
5! \times 2! = 120\times 2 = 240 \text{ formas.}
\]
\end{enumerate}
\end{solucion}

\begin{ejercicio}
Una comisión de 4 personas debe formarse a partir de un grupo de 7 mujeres y 5 hombres (12 personas en total).
\begin{enumerate}[label=(\alph*)]
\item ¿De cuántas formas se puede formar la comisión, sin restricciones?
\item ¿De cuántas formas si la comisión debe tener exactamente 2 mujeres y 2 hombres?
\item Si la comisión se forma completamente al azar, ¿cuál es la probabilidad de que resulte con exactamente 2 mujeres y 2 hombres?
\end{enumerate}
\end{ejercicio}
\begin{solucion}
El orden dentro de la comisión no importa: usamos combinaciones.
\begin{enumerate}[label=(\alph*)]
\item $N=\dbinom{12}{4} = \dfrac{12!}{4!\,8!} = 495$.
\item Elegimos 2 mujeres de las 7 y 2 hombres de los 5:
\[
n_A = \binom{7}{2}\binom{5}{2} = 21\times 10 = 210.
\]
\item Por la regla de Laplace:
\[
P(A) = \frac{n_A}{N} = \frac{210}{495} = \frac{14}{33} \approx 0{,}4242.
\]
\end{enumerate}
\end{solucion}

\begin{ejercicio}
De un mazo de 40 cartas españolas se extraen simultáneamente (sin reposición) 3 cartas. Calcular la probabilidad de que:
\begin{enumerate}[label=(\alph*)]
\item las tres sean ``figuras'' (sota, caballo o rey; hay 12 figuras en total).
\item exactamente 2 sean ases (hay 4 ases en el mazo).
\end{enumerate}
\end{ejercicio}
\begin{solucion}
Casos posibles totales: $N=\dbinom{40}{3} = 9880$.
\begin{enumerate}[label=(\alph*)]
\item Casos favorables: elegir 3 figuras entre las 12 disponibles: $\dbinom{12}{3}=220$.
\[
P(A) = \frac{220}{9880} \approx 0{,}02227.
\]
\item Casos favorables: 2 ases de los 4, y 1 carta no-as de las 36 restantes:
\[
n_B = \binom{4}{2}\binom{36}{1} = 6\times 36 = 216.
\]
\[
P(B) = \frac{216}{9880} \approx 0{,}02186.
\]
\end{enumerate}
\end{solucion}

\begin{ejercicio}
Una fábrica textil produce rollos de tela y, por experiencia histórica, sabe que el $4\%$ de los rollos presenta alguna falla de tejido. Se inspeccionan 5 rollos elegidos al azar de manera independiente entre sí.
\begin{enumerate}[label=(\alph*)]
\item ¿Cuál es la probabilidad de que ninguno de los 5 rollos tenga falla?
\item ¿Cuál es la probabilidad de que al menos uno tenga falla?
\end{enumerate}
\end{ejercicio}
\begin{solucion}
Sea $F_i$: ``el rollo $i$ tiene falla'', con $P(F_i)=0{,}04$ para cada $i$, y los sucesos son independientes entre sí (inspecciones independientes).
\begin{enumerate}[label=(\alph*)]
\item ``Ninguno tiene falla'' $= F_1^c\cap F_2^c\cap F_3^c\cap F_4^c\cap F_5^c$. Por independencia mutua:
\[
P(\text{ninguno con falla}) = (1-0{,}04)^5 = 0{,}96^5 \approx 0{,}8154.
\]
\item El complemento de ``ninguno con falla'' es ``al menos uno con falla'':
\[
P(\text{al menos uno con falla}) = 1-0{,}8154 = 0{,}1846.
\]
\end{enumerate}
\end{solucion}

\begin{ejercicio}
En una urna hay 8 bolillas: 5 blancas y 3 negras. Se extraen 2 bolillas, una a continuación de la otra, \textbf{sin reposición}.
\begin{enumerate}[label=(\alph*)]
\item Calcular la probabilidad de que la segunda bolilla extraída sea blanca, sabiendo que la primera fue blanca.
\item Calcular la probabilidad de que ambas sean blancas (usando la regla del producto).
\item Calcular la probabilidad de que la segunda sea blanca (sin condicionar sobre el resultado de la primera), usando el teorema de la probabilidad total.
\end{enumerate}
\end{ejercicio}
\begin{solucion}
Sean $B_1$: ``la primera es blanca'', $B_2$: ``la segunda es blanca''.
\begin{enumerate}[label=(\alph*)]
\item Si la primera fue blanca, quedan 7 bolillas: 4 blancas y 3 negras.
\[
P(B_2\mid B_1) = \frac{4}{7} \approx 0{,}5714.
\]
\item Por la regla del producto: $P(B_1) = 5/8$, luego
\[
P(B_1\cap B_2) = P(B_1)\cdot P(B_2\mid B_1) = \frac{5}{8}\cdot\frac{4}{7} = \frac{20}{56}=\frac{5}{14}\approx 0{,}3571.
\]
\item $\{B_1,B_1^c\}$ es un sistema completo de sucesos. Necesitamos también $P(B_2\mid B_1^c)$: si la primera fue negra, quedan 5 blancas y 2 negras entre 7, luego $P(B_2\mid B_1^c)=5/7$. Por el teorema de la probabilidad total:
\[
P(B_2) = P(B_1)P(B_2\mid B_1) + P(B_1^c)P(B_2\mid B_1^c) = \frac{5}{8}\cdot\frac{4}{7} + \frac{3}{8}\cdot\frac{5}{7} = \frac{20}{56}+\frac{15}{56}=\frac{35}{56}=\frac{5}{8}.
\]
Notablemente, $P(B_2)=P(B_1)=5/8$: en extracciones sin reposición, la probabilidad marginal de obtener blanca es la misma en cualquier posición de la secuencia.
\end{enumerate}
\end{solucion}

\begin{ejercicio}
Tres máquinas $M_1$, $M_2$ y $M_3$ producen, respectivamente, el $25\%$, el $35\%$ y el $40\%$ del total de piezas de una fábrica. Las proporciones de piezas defectuosas de cada máquina son $5\%$, $4\%$ y $2\%$ respectivamente.
\begin{enumerate}[label=(\alph*)]
\item Calcular la probabilidad de que una pieza elegida al azar del total de la producción sea defectuosa.
\item Si se elige una pieza al azar y resulta defectuosa, calcular la probabilidad de que haya sido producida por la máquina $M_1$.
\item ¿Cuál de las tres máquinas es la ``más sospechosa'' de haber producido una pieza que resultó defectuosa? Justificar con el cálculo de las tres probabilidades a posteriori.
\end{enumerate}
\end{ejercicio}
\begin{solucion}
Sea $D$: ``la pieza es defectuosa''. Datos: $P(M_1)=0{,}25$, $P(M_2)=0{,}35$, $P(M_3)=0{,}40$; $P(D\mid M_1)=0{,}05$, $P(D\mid M_2)=0{,}04$, $P(D\mid M_3)=0{,}02$.
\begin{enumerate}[label=(\alph*)]
\item Por el teorema de la probabilidad total:
\[
P(D) = 0{,}25(0{,}05)+0{,}35(0{,}04)+0{,}40(0{,}02) = 0{,}0125+0{,}014+0{,}008 = 0{,}0345.
\]
\item Por el teorema de Bayes:
\[
P(M_1\mid D) = \frac{P(M_1)P(D\mid M_1)}{P(D)} = \frac{0{,}0125}{0{,}0345} \approx 0{,}3623.
\]
\item Calculamos también las otras dos:
\[
P(M_2\mid D) = \frac{0{,}014}{0{,}0345}\approx 0{,}4058, \qquad P(M_3\mid D) = \frac{0{,}008}{0{,}0345}\approx 0{,}2319.
\]
Verificación: $0{,}3623+0{,}4058+0{,}2319 = 1{,}0000$ (correcto, forman una partición). La máquina \textbf{más sospechosa} es $M_2$, con probabilidad a posteriori $\approx 40{,}6\%$, a pesar de que $M_3$ produce más piezas en total: lo que domina es la combinación entre volumen de producción y tasa de defectos.
\end{enumerate}
\end{solucion}

\begin{ejercicio}
En una universidad, el $60\%$ de los estudiantes de primer año son de la ciudad y el $40\%$ son de otras localidades. La probabilidad de que un estudiante de la ciudad abandone la carrera en el primer año es del $10\%$, mientras que para los estudiantes de otras localidades es del $25\%$.
\begin{enumerate}[label=(\alph*)]
\item Calcular la probabilidad de que un estudiante elegido al azar abandone la carrera en el primer año.
\item Si se sabe que un estudiante \textbf{no} abandonó la carrera, calcular la probabilidad de que sea de otra localidad.
\end{enumerate}
\end{ejercicio}
\begin{solucion}
Sean $C$: ``es de la ciudad'', $C^c$: ``es de otra localidad'', $A$: ``abandona la carrera''. Datos: $P(C)=0{,}60$, $P(C^c)=0{,}40$, $P(A\mid C)=0{,}10$, $P(A\mid C^c)=0{,}25$.
\begin{enumerate}[label=(\alph*)]
\item Probabilidad total:
\[
P(A) = 0{,}60(0{,}10)+0{,}40(0{,}25) = 0{,}06+0{,}10=0{,}16.
\]
\item Necesitamos $P(C^c\mid A^c)$. Primero, $P(A^c)=1-0{,}16=0{,}84$. Además $P(A^c\mid C^c)=1-P(A\mid C^c)=1-0{,}25=0{,}75$. Por Bayes:
\[
P(C^c\mid A^c) = \frac{P(C^c)P(A^c\mid C^c)}{P(A^c)} = \frac{0{,}40\times 0{,}75}{0{,}84} = \frac{0{,}30}{0{,}84}\approx 0{,}3571.
\]
\end{enumerate}
\end{solucion}

\end{document}
